%!TEX root = ./../../springer-mono-format/book.tex
%!TEX root = ./../../thesis-trier.tex

Here, we show the following:
\begin{enumerate}
\item \DFAINT, \DST, and \ST are \LOGSPACE-completeness for partially ordered commutative automata
over a fixed non-empty alphabet (\cref{thm:SF-ST-comm}).
Recall that in \cref{compl:sec:solv} we have shown that these problems
are \NP-complete for partially ordered commutative automata over an unbounded alphabet.

\item \SFS is \LOGSPACE-complete for commutative automata (\cref{thm:SF-ST-comm})
even for unary partially ordered automata. Recall that this problem is \NP-complete
for solvable automata by \cref{thm:SF-solvable}.

\item \DFASWS\ is \LOGSPACE-complete for commutative automata over a fixed alphabet 
and \NP-complete for an unbounded alphabet (\cref{compl:thm:comm-syncshort}).
Hardness remains in both cases even for partially ordered commutative automata.
\end{enumerate}

In this section, we assume an ordering $\Sigma = \{a_1, \ldots, a_k\}$ of each input alphabet~$\Sigma$
with $k \ge 1$ (the case $\Sigma = \emptyset$ is trivial and most problems are solvable in constant time).

\begin{theorem}\label{thm:SF-ST-comm}
	\begin{enumerate}		
		\item \SFS is \LOGSPACE-complete for commutative automata over an unbounded alphabet, where \LOGSPACE-hardness holds even for unary partially ordered automata.
		
		\item \ST, \DST, \DFAINT are \LOGSPACE-complete for partially ordered commutative automata over a fixed non-empty alphabet
		even for an alphabet of size one.
	\end{enumerate}
\end{theorem}
\begin{proof}
	By \cref{compl:prop:reductions},
	where the reductions preserve
	the properties of commutativity, partial orderness,
	and the alphabet,
	it is sufficient to show the second claim for \ST.	
	
	For hardness, we give a reduction from \ONEDAGST that works for \SFS 
	and \ST, and also even if we fix $k$ to be equal to one. 
	For an \ONEDAG-instance $(G, s, t)$ with $G = (V,E)$,
	let $\semiautomat{A}{V}{\Sigma}{\delta}$
	where, for $p \in V$ and $x \in \Sigma$, 
	\[
	\delta(p, x) = \begin{cases}
		q & \text{if } (p,q) \in E \text{ for some } q\in V, \\
		p & \text{otherwise.}		
	\end{cases}
	\]
	
	We have the following:	
	\begin{enumerate}
	\item $\mathcal A$ is a well-defined deterministic and complete automaton
	as for every $p \in V$ there exists at most one $q \in Q$ such that $(p,q) \in E$.
	
	\item $\mathcal A$ is a partially ordered commutative automaton.
		
	\item $t$ is a trap state in $\mathcal A$ since it has no outgoing transitions in~$G$.
	\end{enumerate}

	Now, the following are equivalent:
	\begin{enumerate}
		\item there exists a path from $s$ to $t$ in $G$,
		
		\item $\delta(\{s\}, u) = \{t\}$ for some $u \in \Sigma^*$,
		
		\item $|\delta(\{s,t\}, u)| = 1$ for some $u \in \Sigma^*$,
	\end{enumerate}
	where the latter two properties are equivalent
	as $t$ is a trap state and so must be the synchronizing state.
	By these equivalences,
	the reduction works for the problems \ST and \SFS.
	Thus these problems are \LOGSPACE-hard even for $|\Sigma| = 1$
	by \cref{prelim:thm:complexities:ORD} from \cref{prelim:thm:complexities}.
		
	Next, we show membership in \LOGSPACE.	
	Let $\semiautomat{A}{Q}{\Sigma}{\delta}$ be a partially ordered commutative automaton
	with~$n$ states.
		
	\underline{\SFS}: Let $v = a_1^{n-1}\cdots a_k^{n-1}$. By \cref{prelim:classes:comm:cor:words-SIS-SFS},
	there exists $u \in \Sigma^*$ such that $|\delta(S,u)| = 1$ if and only if
	$|\delta(S,v)| = 1$.
	
	To test whether $|\delta(S,v)| = 1$, we iterate through the states
	$S=\{s_1,\ldots,s_m\}$. For each state $s_i$, we compute $\delta(s_i,v)$
	by processing the letters of $v$ one by one and, for each letter $a_j$,
	applying it repeatedly $n-1$ times. We then check whether all
	computed values $\delta(s_i,v)$ are equal.
	
	This procedure requires three counters: one to iterate over the states of
	$S$, one to iterate over the letters $a_1,\ldots,a_k$, and one to count the
	repetitions $0,\ldots,n-1$ for each letter. In addition, two memory cells
	suffice: one for the current state reached during the computation of
	$\delta(s_i,v)$, and one for storing the result of the previous iteration in
	order to compare $\delta(s_i,v)$ with $\delta(s_{i-1},v)$ for $i>1$.
	
	Using binary encodings for the counters and for the stored states, the total
	space needed is logarithmic.
	
	\underline{\ST}: Let $S,T \subseteq Q$ and suppose, unlike in the proof of \SFS,
	that $\Sigma$ is fixed.
	We iterate through all tuples $(n_1,\ldots,n_k)$ satisfying
	$0 \le n_j < n$ for all $j \in \{1,\ldots,k\}$.
	Each tuple can be stored using $k$ counters, which require
	$k \cdot \log n$ space when encoded in binary.
	
	For each tuple $(n_1,\ldots,n_k)$, let
	$
	w = a_1^{n_1}\cdots a_k^{n_k}.
	$
	We then iterate through all states $q \in S$ and compute $\delta(q,w)$.
	If $\delta(q,w)\in T$ for every $q \in S$, we return \rYES.
	Otherwise, we continue with the next tuple.
	If the tuple $(n-1,\ldots,n-1)$ is reached without finding a 
	word $w$ such that $\delta(q,w)\in T$, we return \rNO.
	
	The computation of $\delta(q,w)$ can be done in logarithmic space.
	Besides the~$k$ counters storing $(n_1,\ldots,n_k)$, we use a counter to
	iterate through the states in $S$, a memory cell storing the current state,
	a counter $j$ going from~$1$ to~$k$ to keep track fo the current letter~$a_j$, and a
	counter $r$ to feed in the factor $a_j^{n_j}$
	of $a_1^{n_1} \cdots a_k^{n_k}$.	
	All of these require only $O(\log n)$ space.
	
	Correctness follows from Equation~\eqref{compl:eqn:com_pto_ST} of
	Corollary~\ref{prelim:classes:cor:com-po}, which implies that whenever there
	exists a word mapping $S$ into $T$, there also exists a word of the form
	$a_1^{n_1}\cdots a_k^{n_k}$ with $0 \le n_j < n$ for all
	$j \in \{1,\ldots,k\}$.
\end{proof}

\begin{theorem}\label{compl:thm:comm-syncshort}
    \DFASWS is \NP-complete even for partially ordered commutative input automata.
	For a fixed non-empty alphabet,
	the problem is \LOGSPACE-complete for commutative input automata  even
	for unary partially ordered automata.
\end{theorem}
\begin{proof}
	\underline{\NP-completeness:} Martyugin~\cite[Prop.~5.3]{Mar2009b}
	proved \NP-hardness for commutative automata over an unbounded alphabet,
	where the reduction resulted in a partially ordered commutative automaton.
	\DFASWS is in \NP\  
	for general input automata. 
	See also the discussion in \cref{compl:sec:subclasses:known-results} (page~\pageref{prelim:classes:intro:DFASWS}).
	
	\medskip\noindent\underline{\LOGSPACE-hardness for fixed unary alphabet:} By \cref{compl:prop:reductions_unary-SWS-SW},
	we have $\DFASWS \eqACZEROm \DFASW$ where the reductions preserve the automaton.
	Then, the claim follows as, by \cref{prelim:thm:finite_list_sync_word}, \DFASW
	is \LOGSPACE-complete for unary partially ordered automata.
		
	\medskip\noindent\underline{\LOGSPACE-membership for fixed alphabet:}
	Let $(\mathcal A, l)$ with $\mathcal A = (Q, \Sigma, \delta)$ 
	be an input instance. By \cref{prelim:classes:comm:prop:set-sync-words:3} from \cref{prelim:classes:comm:prop:set-sync-words},
	we can assume $l < |Q|$.
	Using $k$ counters,
	iterate through all tuples $(n_1, \ldots, n_k)$
	with $n_1 + \cdots + n_k \le l$. By 
	using a binary encoding, this can be done in $O(k \log l)$ space.
	Then, for each tuple, compute 
	$\applytrans{\delta}{q}{a_1^{n_1} \cdots a_k^{n_k}}$ for all $q \in Q$
	and check whether always the same state results.
	If this is the case, then
	return~\rYES. Otherwise, continue with the next tuple. 
	If the iteration through all tuples ends
	and we have not found a synchronizing word, return \rNO.
	
	This algorithm can be made to run in deterministic log-space similarly 
	as in the proof of \cref{thm:SF-ST-comm}.	
\end{proof}

\begin{theorem} \label{prelim:compl:reach-commm}
	\REACH is \NP-complete even for partially ordered commutative automata.
	For a fixed non-empty alphabet,
	the problem is \LOGSPACE-complete even for unary partially ordered automata.
\end{theorem}
\begin{proof}
	Since commutative automata are solvable, 
	by \cref{prelim:compl:thm:reach-NP}, the problem is in \NP for commutative automata.	
	
	For hardness, we give a reduction from \DFASWS
	for partially ordered commutative automata.
	
	Let $(\mathcal A, l)$ with $\mathcal A = (Q, \Sigma, \delta)$
	be an input instance for \DFASWS 
	where $\mathcal A$ is partially ordered and commutative.
	By \cref{prelim:classes:comm:prop:set-sync-words:3} from \cref{prelim:classes:comm:prop:set-sync-words},
	we can assume $l < |Q|$. we can assume $l < |Q| - 1$.
	By \cref{prelim:classes:comm:cor:collections-of-facts:2} if~$\mathcal A$ is synchronizing, 
	the synchronizing
	state is a unique trap state~$t$. Existence of a unique trap state
	and its computation if it exists can be done by an $\AC^0$-circuit
	as in the proof of \cref{compl:prop:reductions:item:st-ss} of \cref{compl:prop:reductions}.
	By \cref{prelim:thm:sync_words_is_ideal}, $\mathcal A$ has a synchronizing
	word of length at most~$l$ if and only if it has one of length exactly~$l$.
	Let $\mathcal A' = (Q', \Sigma, \delta')$
	with $Q' = Q \cup \{q_1, \ldots, q_l, q_{l+1}\}$
	and, for $x \in \Sigma$, we
	set $\delta'(q_i,x) = q_{\min\{l+1,i+1\}}$ 
	and $\delta'(q, x) = \delta(q, x)$ for $q \in Q$.
	Then $\mathcal A$ has a synchronizing word of length~$l$
	if and only if $\delta'(Q', u) = \{t, q_{l}, q_{l+1}\}$
	for some $u \in \Sigma^*$.
	Thus, we obtain a valid reduction to \REACH.
	
	Note that the alphabet is preserved
	and that~$\mathcal A'$ is partially ordered and commutative.
	Thus, by \cref{compl:thm:comm-syncshort},
	for an unbounded alphabet we obtain \NP-hardness
	and for a fixed non-empty alphabet we obtain \LOGSPACE-hardness
	even for partially ordered commutative automata
	and unary partially ordered automata, respectively.\footnote{For \LOGSPACE-hardness
		the reduction from the proof of	 \cref{prelim:prop:finite_list_sync_word-hardness}
		can also be used. In this construction, 
		the set $S = \{ (1, v_j), \ldots, (n, v_j), v_i, (1, v_i), \ldots, (n, v_i) \}$
		is only reachable (by the word $a^{n+j}$ in this case) from the whole state set if $j < i$
		since $(n, v_1) \notin S$
		and so each word reaching~$S$
		has to have length at least~$n$.
		But for $i < j$ this implies $(1, v_j) \notin S$
		since the path  $v_{i+1}$ to $(1, v_j)$
		has length $j - i < n$.}
		
	Next, let $\mathcal A = (Q, \Sigma, \delta)$ be a commutative $n$-state automaton.
	By \cref{prelim:classes:comm:lem:words-SIS-SFS},
	we obtain that all the sets of the form $\delta(Q, u)$ for $u \in \Sigma^*$
	can be written as $\delta(Q, a_1^{i_1} \cdots a_k^{i_k})$
	for $0 \le i_1, \ldots, i_k < n$.	
	
	Now, the subsets $\delta(Q, a_1^{i_1} \cdots a_k^{i_k})$
	for $0 \le i_1, \ldots, i_k < n$	
	can be iterated through and tested for equality
	with a given target set $S$ in deterministic log-space
	similarly as in the proof of Theorem~\ref{thm:SF-ST-comm} for \ST 
	by iterating through the states, the letters, the numbers $i_j$, 
	and using two memory cells. We also use, additionally, $|S|$ Boolean flags
	that indicate for each state in~$S$ whether it has
	been hit by some state when applying $a_1^{i_1} \cdots a_k^{i_k}$
	to the state set. Then the algorithm
	returns \rYES precisely if all these flags are set to true
	for some $0 \le i_1, \ldots, i_k < n$.
	For a fixed alphabet, 
	this gives a deterministic algorithm using only logarithmic space.
\end{proof}

\begin{theorem}  \label{prelim:compl:reachf-commm}
	\REACHF is \LOGSPACE-complete for  partially ordered commutative automata
	over a fixed non-empty alphabet, even for an alphabet of size one.
\end{theorem}
\begin{proof}
	By \crefitem{compl:prop:reductions}{compl:prop:reductions:item:REACH-REACHF},
	we have $\REACH \leACZEROm \REACHF$ and the reduction
	preserved the automaton type and the alphabet.
	This yields \LOGSPACE-hardness by \cref{prelim:compl:reach-commm}.
	
	Next, the problem on commutative partially ordered automata can be solved 
	with the same algorithm given in the proof
	of Theorem~\ref{thm:SF-ST-comm} for \ST, but where we also have
	to check for each tuple $(n_1, \ldots, n_k)$
	not only that $\delta(q, a_1^{n_1}\cdots a_k^{n_k})$
	is in~$T$, but that each state from~$T$
	appears at least once for some~$q \in S$, i.e., the image equals precisely~$T$.
\end{proof}

\begin{theorem} \label{prelim:compl:thm:reach-comm-fixed}
	\REACHF is solvable in polynomial time for commutative automata 
	when the alphabet is fixed.	
\end{theorem}
\begin{proof}
	In~\cite[Thm.~1.4]{DBLP:conf/soda/BabaiBCIL96}, for fixed~$h$ and~$r$,
	commuting $n \times n$ matrices $A_1, \ldots, A_h$,
	$B_1,\ldots, B_r$, and $n \times l$ matrices $A$ and $B$,
	it was proved that it is decidable in polynomial time whether
	there exist non-negative integers $x_i$ and $y_j$ such that
	\begin{equation} \label{prelim:compl:eqn:reach-comm-fixed}
		\left( \prod_{i=1}^h A_i^{x_i} \right)A 
		= 
		\left( \prod_{j=1}^r B_j^{y_j} \right)B.
	\end{equation}
	Suppose a commutative automaton $\semiautomat{A}{Q}{\Sigma}{\delta}$
	with $\Sigma = \{a_1, \ldots, a_h\}$
	and $S, T \subseteq Q$ is given.
	Let $A_i$ be the transition matrix of~$a_i$,
	let $r = 1$ and $B_1$ be the identity matrix,
	and let $A$ and $B$ be the $n \times 1$ matrices
	resulting from the characteristic vectors
	of $S$ and~$T$, respectively.
	Then  \cref{prelim:compl:eqn:reach-comm-fixed}
	with these matrices has a solution $x_1, \ldots, x_h \ge 0$ 
	if and only if $\delta(S, a_1^{x_1} \cdots a_h^{x_h}) = T$.
	Hence, 
	we find that \REACH 
	is solvable in polynomial time for fixed~$h$
	and commutative~$\mathcal A$.
\end{proof}